\documentclass{article}
\usepackage{graphicx}
\usepackage[slovene]{babel}
\usepackage{amsfonts,amssymb,amsmath, mathtools}
\usepackage[utf8]{inputenc}
\usepackage[T1]{fontenc}
\usepackage{lmodern}
\usepackage{graphicx} % Required for inserting images
\usepackage{dsfont}  % za font za nekatere oznake

% \usepackage{semantic}  % za inference rules, type theory
\usepackage{proof}


% \usepackage{mathtools}  % https://tex.stackexchange.com/questions/45001/how-do-i-left-align-entries-in-a-matrix-with-beginmatrix
\usepackage{amsthm}  % za theoremstyle ukaz

\PassOptionsToPackage{hyphens}{url}\usepackage{hyperref} % https://tex.stackexchange.com/questions/3033/forcing-linebreaks-in-url
% dejansko to ni le za linke do spletnih strani, ampak tut za povezave do slik/zgledov ...

\DeclareMathOperator{\Sym}{Sym}
\DeclareMathOperator{\Perm}{Perm}
\DeclareMathOperator{\supp}{supp}
\DeclareMathOperator{\Terms}{Terms}
\DeclareMathOperator{\sub}{sub}

\newcommand{\inference}[3][]{\infer[#1]{#3}{#2}}

\def\N{\mathbb{N}} % mnozica naravnih stevil
\def\Z{\mathbb{Z}} % mnozica celih stevil
\def\Q{\mathbb{Q}} % mnozica racionalnih stevil
\def\R{\mathbb{R}} % mnozica realnih stevil
\def\C{\mathbb{C}} % mnozica kompleksnih stevil
\newcommand{\geslo}[2]{\noindent\textbf{#1} \quad \hangindent=1cm #2\\[-1pc]}

\def\qed{$\hfill\Box$}   % konec dokaza
\def\qedm{\qquad\Box}   % konec dokaza v matematičnem načinu
\newtheorem{izrek}{Izrek}
\newtheorem{trditev}{Trditev}
\newtheorem{posledica}{Posledica}
\newtheorem{lema}{Lema}
\newtheorem{opomba}{Opomba}
\newtheorem{definicija}{Definicija}
%\theoremstyle{remark}
\newtheorem{zgled}{Zgled}
%\theoremstyle{plain}
\newtheorem{domneva}{Domneva}
\newtheorem{dokaz}{Dokaz}  % verjetno neuporabno
\newtheorem{oznaka}{Oznaka}
\newtheorem{izrek_crosser}{Izrek (Church–Rosserjev izrek)}


\title{Lambda izrazi v nominalnih množicah}
\author{Tadej Glinšek}
\date{Maj 2026}





\begin{document}

\maketitle

% \section{Introduction}


%\section{Nominal Sets}

\section{Predznanje}
Najprej uvedimo potrebno predznanje, ki ga bomo potrebovali pri obravnavi. 

V poglavju bomo z \(G\) označevali grupo, ki je definirana kot množica z binarno operacijo, ki je asociativna, ima enoto in je inverzibilna. Spomnimo se, kaj pomeni delovanje grupe na množico.
\begin{definicija}
    
Množica \(X\) je \emph{\(G\)-množica}, če grupa \(G\) deluje na množici \(X\). Delovanje je definirano na sledeč način:
\[
    (g_1 \cdot g_2) \cdot x = g_1 \cdot (g_2 \cdot x)
\]
\[
    e \cdot x = x,
\]
kjer je \(e\) enota grupe \(G\), \(g_1,\ g_2 \in G\) dva elementa grupe in \(x\) element množice \(X\).
\end{definicija}

Predvsem bo za nas uporabno delovanje simetrične grupe množice na množico samo.

\begin{definicija}
    Naj bo \(X\) množica. Potem s \(\Sym(X)\) označimo njeno simetrično grupo, tj. grupo vseh permutacij te množice.
\end{definicija}

Množenje permutacij v simetrični grupi je kar kompozitum permutacij. Delovanje simetrične grupe torej izgleda takole.

\begin{zgled}
Elementa \(\pi_1,\ \pi_2 \in \Sym(X)\) delujeta na \(X\) takole:
\[\pi_1 \cdot \pi_2 \cdot x = \pi_1 \cdot (\pi_2(x)) = \pi_1(\pi_2(x)) = (\pi_1 \circ \pi_2)(x) = (\pi_1 \circ \pi_2) \cdot x\]
\end{zgled}

Za nas bo množica \(X\) največkrat neskončna. Ampak večinoma nas bodo zanimale permutacije, ki permutirajo le končno mnogo elementov te množice. Zato uvedimo naslednjo definicijo.

\begin{definicija}
Permutacija \(\pi \in \Sym(A)\) je \textbf{končna}, če je \(\{a \in A\ |\ \pi a \neq a\}\) končna množica.
\end{definicija}

Množico končnih permutacij iz \(\Sym(A)\) označimo s \(\Perm(A)\). To je tudi grupa s podedovano operacijo kompozituma permutacij.


\section{Nominalne množice}

V tem poglavju bomo z \(\mathds{A}\) označili števno neskončno množico elementov, ki jim pravimo \emph{atomska imena} (oz. \textbf{atomi}).
Prav tako bo v tem poglavju \(X\) vedno \(\Perm \mathds{A}\)-množica.

\textbf{Nosilec} (angl. \emph{support}) elementa \(x \in X\) definiramo kot množico \(A \subseteq \mathds{A}\), za katero velja, da za vsak \(\pi \in \Perm \mathds{A}\) velja, da če \(\pi\) fiksira vsak \(a \in A\), potem \(\pi\) stabilizira \(x\). Z drugimi besedami, \(A\) je nosilec \(x\)-a, če:
\[
    \forall\ \pi \in \Perm \mathds{A}.\ (\forall\ a \in A.\ \pi\ a = a) \Rightarrow \pi \cdot x = x
\]

Nosilec elementa \(x \in X\) označimo \(\supp_X(x)\). Če lahko množico \(X\) razberemo iz konteksta, potem pišemo kar \(\supp(x)\).

Zakaj smo za množico elementov, za katere velja, da če so ti fiksirani je fiksiran tudi \(x \in X\), izbrali ime \emph{nosilec}? Razlog je v tem, da kako se ti elementi obnašajo, to določa, kaj se zgodi z \(x\). Elementi so tako kot ogrodje: dokler se ogrodje ne spremeni, se tudi \(x\) ne more spremeniti. Nosilec torej določa vse kar je pomembno za \(x\). V tem smislu se pojem nosilca ne razlikuje preveč od pojma nosilca pri funkcijah, kjer nosilec sestavlja podmnožico domene, kjer vrednost funkcije ni ničelna.


Obstaja naslednja karakterizacija nosilca:

\begin{trditev}
    Množica \(A \subseteq \mathds{A}\) je nosilec za \(x \in X\) natanko tedaj, ko 
    \[
        \forall\ a,\ a' \in \mathds{A}\setminus A.\ (a\ a') \cdot x = x
    \]
\end{trditev}

\noindent
{\em Dokaz:\/} Zaenkrat dokaz opustimo, povejmo samo idejo:

Če transpozicije stabilizirajo \(x\), potem tudi produkt le-teh stabilizira \(x\). Vsaka končna permutacija se lahko zapiše kot produkt transpozicij. Permutacija, ki je produkt transpozicij, ki fiksirajo elemente iz \(A\), bo prav tako fiksirala elemente iz \(A\).    \qed

\begin{opomba}
    Nosilec elementa \(x\) si lahko predstavljamo kot množico vseh \(a\), ki delujejo na \(x\) na netrivialen način (tj., niso v njegovem stabilizatorju).
\end{opomba}

\begin{zgled}
    Vsak \(B \supseteq A\), kjer je \(A\) nosilec elementa \(x\), je spet nosilec elementa \(x\).
\end{zgled}

\begin{definicija}
    Če za element \(x \in X\) obstaja nosilec \(A\), ki je končna množica, potem je \(x\) \textbf{končno podprt} (angl. \emph{finitely supported}).  % podprt? nošen?
\end{definicija}

Predpostavimo, da ima \(x\) dva nosilca \(A\) in \(A'\), tako da \(A\, \cup\, A'\) ne pokrije celotnega \(\mathds{A}\) (tj. \(\mathds{A} \setminus (A\, \cup\, A')\) je neprazen). Potem se izkaže, da je \(A\, \cap\, A'\) prav tako nosilec \(x\).

\noindent
{\em Dokaz:\/}
Dokazujemo
    \[
        \forall\ a,\ a' \in \mathds{A}\setminus (A\, \cap\, A').\ (a\ a') \cdot x = x
    \]

Če sta \(a\) in \(a'\) oba iz \(\mathds{A}\setminus A\), potem \((a \ a') \cdot x = x\), saj je \(A\) nosilec. Analogno sklepamo, če sta \(a\) in \(a'\) oba iz \(\mathds{A}\setminus A'\).  % dejansko tega ne rabimo, saj argument deluje v splošnem
Sicer predpostavimo (BŠS) \(a \in \mathds{A}\setminus A\) in \(a' \in \mathds{A}\setminus A'\), velja pa naj tudi \(a,\ a' \notin (\mathds{A} \setminus A)\, \cap\, (\mathds{A} \setminus A')\). V posebnem velja \(a \neq a'\).

Velja \(\mathds{A}\setminus (A\, \cap\, A') = (\mathds{A} \setminus A)\, \cup\, (\mathds{A} \setminus A')\).
Prav tako \(\mathds{A} \setminus (A\, \cup\, A') = (\mathds{A} \setminus A)\, \cap\, (\mathds{A} \setminus A') \), kar pa po predpostavki ni prazna množica.
Izberimo \(a'' \in (\mathds{A} \setminus A)\, \cap\, (\mathds{A} \setminus A')\).
Velja \(a \neq a'' \neq a'\), namreč če bi katerakoli izmed enakosti veljala, potem bi vsaj eden izmed \(a\) in \(a'\) bil v \((\mathds{A} \setminus A)\, \cap\, (\mathds{A} \setminus A')\), kar pa je v protislovju s predpostavko.

Ker sta \(a,\ a''\) oba iz \(\mathds{A} \setminus A\), velja \((a \ a'') \cdot x = x\). Podobno, ker sta \(a',\ a''\) oba iz \(\mathds{A} \setminus A'\), velja \((a' \ a'') \cdot x = x\).

S pomočjo \(a''\) lahko \(a\) in \(a'\) zamenjamo na sledeč način:
\[(a \ a') = (a \ a'') \circ (a' \ a'') \circ (a \ a'')\]

Torej velja \((a \ a') \cdot x = x\). \qed

\begin{posledica}
    Če sta \(A \subseteq \mathds{A}\) in \(A' \subseteq \mathds{A}\) končna nosilca elementa \(x\), je tudi \(A\, \cap\, A'\) nosilec elementa \(x\).
\end{posledica}

\begin{definicija}
    \textbf{Nominalna množica} je \(\Perm\mathds{A}\)-množica, katere vsak element ima končen nosilec.
\end{definicija}

Delovanje grupe \(\Perm\mathds{A}\) na množici je lahko poljubno; odvisno, kako ga definiramo.

Alternativno lahko na nosilec gledamo na sledeč način. 
Nominalni množici \(X\) lahko priredimo preslikavo \(\supp: X \rightarrow \mathcal{P}_{\text{fin}} \mathds{A}\), kjer je \(\mathcal{P}_{\text{fin}} \mathds{A}\) množica vseh končnih podmnožic \(\mathds{A}\).

\subsubsection{Produkti nominalnih množic}

Vprašajmo se, ali je kartezični produkt nominalnih množic spet nominalna množica. Izkaže se, da je to res, če imamo končni produkt - takoj ko množimo neskončno mnogo nominalnih množic, se lahko stvari zapletejo. Ampak za nas bo zares relevanten le končni produkt.

Naj bodo \(X_i\), za \(i = 1,\ \ldots,\ n\), nominalne množice. Produkt \(X_1 \times \cdots \times X_n\) je potem nominalna množica: če za \(\pi \in \Perm\mathds{A}\) in \((x_1,\ \ldots,\ x_n) \in X_1 \times \cdots \times X_n\) definiramo delovanje 
    \[
        \pi \cdot (x_1,\ \ldots,\ x_n) = (\pi \cdot x_1,\ \ldots,\ \pi \cdot x_n)
    \]
potem se ni težko prepričati, da je nosilec tega elementa spet končen. Specifično, za nosilec \(n\)-terice velja
\[
\supp_{X_1 \times \cdots \times X_n}(x_1,\ \ldots,\ x_n) = \supp_{X_1} x_1\, \cup\, \cdots\, \cup\, \supp_{X_n} x_n.
\]

\subsection{Relacije nad nominalnimi množicami}

Struktura nominalne množice \(X\) je določena s tem, kako definiramo delovanje grupe \(\Perm\mathds{A}\) na tej množici. Toda načeloma nas ne zanimajo specifične permutacije, ampak predvsem kako grupa \(\Perm\mathds{A}\) kot celota deluje na nekem specifičnem elementu. Primer take definicije je nosilec, ki nam za nek element pove, kateri elementi določijo delovanje tega elementa.

Zato uvedimo še naslednje relacije na elementih \(X\), kasneje pa bomo dokazali lastnosti, s katerimi bomo lahko logično sklepali, brez da se sklicujemo na delovanje sámo.

Vpeljimo najprej relacijo orbitne ekvivalence:
\begin{definicija}
    Naj bosta \(x, y\) elementa nominalne množice \(X\). Rečemo, da sta \(x\) in \(y\) \textbf{orbitno ekvivalentna}, če velja
    \[
        x \sim y :\Longleftrightarrow \exists\ \pi \in \Perm \mathds{A}. \pi \cdot x = y.
    \]
\end{definicija}
Z drugimi besedami, velja \(x \sim y\), če sta oba v isti orbiti delovanja grupe \(\Perm \mathds{A}\) na \(X\). Izkaže se, da je ta relacija je tudi nominalna množica kot množica parov elementov, ki so si v relaciji.

Obstaja še pogojni analog te relacije:
\begin{definicija}
    Naj bosta \(x, y\) elementa nominalne množice \(X\) in \(z\) element nominalne množice \(Z\). Rečemo, da sta \(x\) in \(y\) \textbf{orbitno ekvivalentna glede na \(z\)}, če velja
    \[
        x \sim y \ \vert \ z :\Longleftrightarrow \exists\ \pi \in \Perm \mathds{A}. \pi \cdot x = y\, \wedge\, \pi \cdot z = z.
    \]
\end{definicija}

Naslednji relaciji, ki ju uvajamo, se poslužujeta nosilcev elementov.
\begin{definicija}
    Naj bosta \(X\) in \(Y\) nominalni množici in \(x \in X\), \(y \in Y\). Če za \(x\) in \(y\) velja
    \[
        x \ \#\ y :\Longleftrightarrow \supp x\, \cap\, \supp y = \emptyset,
    \]
    potem rečemo, da sta \textbf{separirana}, relaciji pa \textbf{separiranost}.
    
    Če je \(Z\) še ena nominalna množica in je \(z \in Z\) elementov, potem če velja
    \[
        x \ \#\  y \ \vert \ z :\Longleftrightarrow \supp(x)\, \cap\, \supp(y) \subseteq \supp(z)
    \]
    rečemo, da sta \(x\) in \(y\) \textbf{separirana glede na} \(z\). Relaciji pravimo \textbf{pogojna separiranostna relacija}.  % ali odvisna?
\end{definicija}
Ta druga definicija je pogojni analog prve definicije.
Izkaže se, da sta obe relaciji spet nominalni množici sami po sebi.
\subsubsection{Lastnosti relacij}

Poglejmo si nekaj lastnosti, ki veljajo za relacijo orbitne ekvivalence.

\begin{trditev}
    Naj bosta \(X\) in \(Y \subseteq X\) nominalni množici. Naj \(x \in X\) in \(y \in Y\). Potem \(x \sim y \Rightarrow x \in Y\).
\end{trditev}

Ta trditev nam pove, da so podnominalne množice unije ekvivalenčnih parov relacije \(\sim\).





Naslednja trditev združuje relacijo orbitne ekvivalence in relacijo separiranosti.

\begin{trditev}
    Za vsaka \(x,\ y \in X\) obstaja \(z \in X\), da velja \(z \ \#\  x\) in \(z \sim y\).
\end{trditev}

Sledi še pogojna oblika prejšnje trditve.

\begin{trditev}
    \label{pogojna_trditev}
    Naj bodo \(X, Y, W\) nominalne množice in \(x \in X\), \(y \in Y\), \(w \in W\) njihovi elementi. Potem obstaja \(z \in X\), da velja \(z \sim x\ \ \vert\ \ w\) in \(z \ \#\  y\ \ \vert\ \ w\).
\end{trditev}


Dokaze zaenkrat opustimo.

Naslednja trditev nam bo prišla prav:
\begin{trditev}

    Naj bodo \(X,\ Y,\ Z,\ W\) nominalne množice, \(x,\ y,\ z,\ w\) pa njihovi elementi. Če sta \(a \in A\) in \(b \in B\) elementa nominalnih množic, potem z izrazom \((a,\ b)\) označujemo element produkta \(A \times B\). Potem veljajo naslednje implikacije:
    \begin{equation}
        \label{a1}
        x \ \#\  (y,\ z) \ \vert\ w \Longleftrightarrow x \ \#\  y \ \vert\ (z,\ w),
    \end{equation}
    \begin{equation}
        \label{a2}
        x \ \#\  (y,\ z) \ \vert\ w \Longleftrightarrow x \ \#\  z \ \vert\ w,
    \end{equation}

    in
    \begin{equation}
        \label{a3}
        x \ \#\  y \ \vert\ (z,\ w)\, \wedge\, x \ \#\ z \ \vert\ w \Longleftrightarrow x \ \#\ (y,\ z) \ \vert\ w.
    \end{equation}
\end{trditev}
Te trditve lahko tudi skupaj zapišemo kot
\[
    x \ \#\  (y,\ z) \ \vert\ w \Longleftrightarrow x \ \#\  y \ \vert\ (z,\ w)\, \wedge\, x \ \#\  z \ \vert\  w.
\]
\newpage
\section{Netipiziran lambda račun}

\iffalse
Množico izrazov v netipiziranem lambda računu označujemo z \(\Lambda\). Definirani so s pravili:
\[
    t \in \Lambda ::= a \ | \ \lambda a.t\ | \ t\ t
\]
kjer je \(a \in \mathds{A}\).

Ideja je, da izrazov ne ločujemo tako strogo, kot jih ločuje običajna enakost. Namreč, \(\Perm \mathds{A}\) lahko na izraze deluje, tako da med sabo permutira elemente \(\mathds{A}\) (atome), ki se v izrazih pojavijo. Permutiranje atomov ne spremeni vsebine izrazov, zato tudi take izraze ne bomo ločevali.

Z \(\Gamma\) bomo označevali končno zaporedje atomov iz \(\mathds{A}\). To predstavlja kontekst prostih spremenljivk, ki se pojavi v naših izrazih (angl. terms). Torej izraze bomo pisali z eksplicitnim kontekstom, npr. \(\Gamma,\ y,\ z \vdash \lambda x.xy\)

Naj bo \(\sim\) ekvivalenčna relacija, ki enači izraze oblike \(\Gamma \vdash M\), ki so v isti orbiti glede na delovanje grupe \(\Perm{\mathds{A}}\). Informalno, ko permutacija deluje na izraz, permutira spremenljivke, ki se pojavijo v tem izrazu in kontekstu, kjer izraz živi. Kontekst je namreč zaporedje elementov množice \(\mathds{A}\). 

Sedaj definiramo relacijo \(\equiv_{\alpha}^\Gamma\), ki je relacija med izrazi nad istim kontekstom. Natančneje, ta relacija je relacija na množici 
\[
    \Terms_{\Gamma} \coloneq \{\Gamma \vdash M |\ \Gamma \vdash M \ \text{je izraz}\}.
\]
To je množica izrazov nad istim kontekstom \(\Gamma\).

Tej ekvivalenci pravimo \textbf{\(\alpha\)-ekvivalenca}. Namesto \(\Lambda\) se bomo torej pogosto osredotočili na \(\Lambda/\equiv_{\alpha}\).

Množica vseh izrazov je nominalna množica:

\[
    \Terms \coloneq \{\Gamma \vdash M |\ \Gamma \vdash M \ \text{je izraz}\}
\]

\fi


V tem poglavju bomo definirali lambda izraze na drugačen način kot smo jih vajeni. Namreč, v definicijo lambda izraza bomo vgradili kontekst, ki bo hranil vse možne proste spremenljivke izraza.

\begin{definicija}
    \textbf{Kontekst} \(\Gamma\) je končno zaporedje paroma različnih atomov.        
\end{definicija}


Iz \(\lambda\)-računa vemo, kako so definirani izrazi \(M\), če predpostavimo nek konstekst \(\Gamma\). Mi se bomo tega lotili drugače: \(\lambda\)-izraze bomo definirali skupaj s kontekstom.

\begin{definicija}
    \(\lambda\)-izraz v kontekstu \(\Gamma \vdash M\) je definiran z naslednjimi pravili:
    \[
        \inference[(x \in \Gamma)]{}{\Gamma \vdash x}
    \]
    \[
        \inference{\Gamma \vdash M \quad \Gamma \vdash N}{\Gamma \vdash M N}
    \]
    \[
        \inference[(x \notin \Gamma)]{\Gamma,\ x \vdash M}{\Gamma \vdash \lambda x. M}    
    \]
\end{definicija}

Kontekst \(\Gamma\) predstavlja nabor možnih prostih spremenljivk, ki se lahko pojavijo v našem lambda izrazu. Torej, izraze bomo pisali z eksplicitnim kontekstom, npr. \(\Gamma,\ y,\ z \vdash \lambda x.xy\)

\begin{opomba}
    Če se \(x\) pojavi v \(\Gamma \vdash M\), potem se ne more pojaviti kot prosta in vezana spremenljivka hkrati.
\end{opomba}



\subsection{Delovanje permutacij na množici \(\Terms\)}

Zdaj moramo definirati delovanje permutacij na lambda izrazih. Najprej definiramo, kako permutacija deluje na kontekstu:

\begin{definicija}
    Naj bo \(\Gamma = a_1,\,\ a_2,\ \ldots,\ a_n\) končno zaporedje elementov iz \(\mathds{A}\). Permutacija \(\pi \in \Perm\mathds{A}\) slika zaporedje \(\Gamma\) v \(\pi(a_1), \ \pi(a_2), \ \ldots, \ \pi(a_n)\).
\end{definicija}

Definicija pravi, da permutacija deluje na kontekstu po atomih.

\begin{opomba}
    Množica vseh kontekstov \(\Gamma\) je nominalna množica.
\end{opomba}


Sedaj lahko definiramo delovanje permutacije na lambda izrazu nad kontekstom. Ideja je, da permutiramo lambda izraz in njegov kontekst \emph{po atomih}:

\begin{definicija}
    Delovanje grupe \(\Perm{\mathds{A}}\) na množico izrazov je definirano sledeče:
    \[
        \pi(\Gamma \vdash x) \coloneq \pi(\Gamma) \vdash \pi(x)
    \]
    \[
        \pi(\Gamma \vdash MM') \coloneq \pi(\Gamma ) \vdash \pi(M)\pi(M')
    \]
    \[
        \pi(\Gamma \vdash \lambda x.\ M) \coloneq \pi(\Gamma) \vdash \lambda \pi(x). \pi(M)
    \]
\end{definicija}

\begin{opomba}
    Izraz \(\pi(\Gamma) \vdash \lambda \pi(x). \pi(M)\) je dobro definiran, saj se \(\pi(x)\) ne more pojaviti v \(\pi(\Gamma)\), če se \(x\) ne pojavi v \(\Gamma\).
\end{opomba}

Ker smo s tem opisali delovanje permutacij na množico izrazov oblike \(\Gamma \vdash M\), je ta množica \(\Perm{\mathds{A}}\)-množica. Izkaže se, da je tudi nominalna množica.

Lambda izraze nad kontekstom lahko gledamo kot množico:
\begin{definicija}
    \[
        \Terms_{\Gamma} \coloneq \{\Gamma \vdash M\ |\ \Gamma \vdash M \ \text{je \(\lambda\)-izraz}\}
    \]
    \[
        \Terms \coloneq \{\Gamma \vdash M\ |\ \Gamma\ \text{je kontekst},\ \Gamma \vdash M \ \text{je \(\lambda\)-izraz}\}
    \]
\end{definicija}

Obedve množici sta nominalni množici.

\begin{opomba}
    Včasih bomo kontekst spustili. Takrat \(\pi(M)\) pač predstavlja izraz \(M'\) v \(\pi(\Gamma \vdash M) = \Gamma' \vdash M'\).
\end{opomba}

\begin{opomba}
    Velja \(\pi(\supp(M)) = \supp(\pi(M))\). To sicer velja za splošne nominalne množice.
\end{opomba}

Še opomba glede oznak.
Spomnimo se relacije orbitne ekvivalence: za dva izraza \(\Gamma \vdash M\) in \(\Gamma' \vdash N\) velja \(\Gamma \vdash M \sim \Gamma' \vdash N\) (pišemo tudi \(M,\ \Gamma \sim N,\ \Gamma'\)), kadar obstaja \(\pi \in \Perm{\mathds{A}}\), da \(\pi(\Gamma \vdash M) = \Gamma' \vdash N\) oz. alternativno \(\pi(M,\ \Gamma) = N,\ \Gamma'\). Če velja \(\Gamma = \Gamma'\), potem lahko pišemo \(\Gamma \vdash M \sim N\) ali pa \(M \sim N \ \vert\ \Gamma\).

Opomba: \(\pi(\Gamma) = \Gamma'\) pomeni, da sta obe zaporedji iste dolžine in da se \(i\)-ti člen zaporedja \(\Gamma\) slika v \(i\)-ti člen zaporedja \(\Gamma'\). Torej, \(\Gamma \sim \Gamma'\) pomeni, da za vsako mesto v zaporedju imamo isto permutacijo \(\pi\): ne more \(\pi\) biti odvisna od \(i\), kjer \(i\) pomeni indeks člena v zaporedju. To je sorodno s konceptom \emph{odvisne porazdelitve} (angl. joint distribution) pri slučajnih spremenljivkah.

\subsection{Relacije lambda izrazov}

Zdaj, ko imamo nominalno strukturo na lambda izrazih, lahko pogledamo, kaj te relacije pomenijo.

Po definiciji delovanja permutacije na lambda izrazu, velja, da sta dva izraza orbitno ekvivalentna, če lahko s permutacijo spremenljivk (tako prostih kot vezanih) pridemo iz enega izraza v drugega. Ampak treba je paziti, da ne spreminjamo vrstnega reda spremenljivk v kontekstu: če spremenimo, načeloma dobimo izraz, ki ni več orbitno ekvivalenten prvotnemu.

Velikokrat bomo imeli dva izraza s podobnim kontekstom, kjer želimo, da se njima skupne spremenljivke ne permutirajo, vse preostale spremenljivke se pa lahko. Zato lahko uvedemo naslednjo notacijo.

\begin{oznaka}
    Če imamo izraza \(\Gamma \vdash M\) in \(\Delta \vdash N\), potem namesto \((\Gamma \vdash M) \sim (\Delta \vdash N) \ \vert\ (\Gamma\, \cap\, \Delta)\) pišemo kar \((\Gamma\, \cap\, \Delta) \vdash M \sim N\). Pogosto namreč fiksiramo elemente obeh kontekstov, ki so skupni obema izrazoma.
\end{oznaka}

Zapis je malce dvoumen, ker sta \(\Gamma\) in \(\Delta\) načeloma večja konteksta kot pa \((\Gamma\, \cap\, \Delta)\). Toda načeloma je razvidno, kaj sta konteksta od \(M\) in \(N\), in ker se pojavita le v enem izmed obeh izrazov, nista tako relevantna za razumevanje pomena te relacije.

Povsem analogno naredimo z relacijo separiranosti:

\begin{oznaka}
    Če imamo izraza \(\Gamma \vdash M\) in \(\Delta \vdash N\), potem namesto \((\Gamma \vdash M) \ \#\ (\Delta \vdash N) \ \vert\ (\Gamma\, \cap\, \Delta)\) pišemo kar \((\Gamma\, \cap\, \Delta) \vdash M \ \#\ N\). Pogosto namreč fiksiramo elemente obeh kontekstov, ki so skupni obema izrazoma.
\end{oznaka}


Poglejmo si primer, kako lahko teorijo nominalnih množic uporabimo pri relacijah med lambda izrazi:

\begin{lema}
    Naj bo \(\Gamma \vdash M\) izraz nad kontekstom \(\Gamma\) in naj bo \(\Delta \vdash N\) izraz nad kontekstom \(\Delta\), tako da velja \((\Gamma \vdash M) \sim (\Delta \vdash N) \ \vert\ \Gamma\). Potem je \(\Gamma \vdash N\) dobro definiran izraz.
\end{lema}

\noindent
{\em Dokaz:\/}
Dokaz zaenkrat opustimo.
\qed

\subsection{Substitucija lambda izrazov}

Uvedimo pojem substitucije.

\begin{definicija}
    Naj bosta \(\Gamma,\ x,\ \Gamma' \vdash M\) in \(\Gamma,\ \Gamma' \vdash N\) \(\lambda\)-izraza, tako da velja \[N \ \#\  M \ \vert\  \Gamma,\ \Gamma'.\] Potem definiramo \(\Gamma,\ \Gamma' \vdash M[N/x]\) kot
    
    \[z[N/x]\, \coloneq\, \begin{cases*}
    N & če z = x,\\
    z & sicer,
    \end{cases*}
    \]
    \[
        (M'\ M'')[N/x]\, \coloneq\, M'[N/x]\ M''[N/x],\ \ \ \ \ 
    \]
    \[
        (\lambda z.\ M')[N/x]\, \coloneq\, \lambda z.\ (M'[N/x]).\quad \quad \quad
    \]
    
    
\end{definicija}

\begin{izrek}
    Izraz \(M[N/x]\) je izraz induciran s kontekstom \(\Gamma,\ \Gamma'\), tj., \(\Gamma,\ \Gamma' \vdash M[N/x]\).
\end{izrek}

Izrek pravi, da je ob določenih predpostavkah izraz \(\Gamma \vdash M[N/x]\) dobro definiran \(\lambda\)-izraz. Dokaz ne bo tako očiten, kot se morda zdi na prvi pogled.

Za dokaz bomo potrebovali naslednjo lemo:

\begin{lema}
    Naj \(\Gamma \vdash M\) \(\lambda\)-izraz, \(z\) spremenljivka, in naj velja \(z \ \#\  M,\ \Gamma\). Potem velja \[
        \Gamma,\ z \vdash M.
    \]
\end{lema}

\noindent
{\em Dokaz:\/}

Naj \(M\) izraz v kontekstu \(\Gamma\). Naj \(z\) spremenljivka, tako da \(z \ \#\  M\) in \(z \notin \Gamma\).

\begin{itemize}
    \item Recimo, da \(M = x\), kjer je \(x\) neka spremenljivka. Velja \(x \neq z\), saj velja \(z \ \#\  M\). Iz \(\Gamma \vdash M\) sledi \(x \in \Gamma\), torej \(x \in \Gamma,\ z\). Posledično, skupaj z dejstvom da \(z \notin \Gamma\), velja \(\Gamma,\ z \vdash M\).

    \item Recimo, da \(M = M'M''\). Sledi \(\Gamma \vdash M'\) in \(\Gamma \vdash M''\). Iz \(z \ \#\  M\) sledi \(z \ \#\  M'\) in \(z \ \#\  M''\). Potem so izpolnjeni pogoji za trditev za \(M'\) in \(M''\), torej velja \(\Gamma,\ z \vdash M'\) in \(\Gamma,\ z \vdash M''\). Torej \(\Gamma,\ z \vdash M'M''\), kar smo želeli dokazati.

    \item Recimo, da \(M = \lambda x. M'\). Potem \(\Gamma,\ x \vdash M'\). Iz \(z \ \#\  \lambda x. M'\) sledi \(z \ \#\  M'\) in \(x \neq z\). Torej lahko uporabimo indukcijsko predpostavko in dobimo \(\Gamma,\ x,\ z \vdash M'\). Sledi \(\Gamma,\ z \vdash \lambda x. M'\).
\end{itemize}

\qed

Sedaj lahko dokažemo izrek:

\noindent
{\em Dokaz izreka:\/}

Predpostavimo \(\Gamma,\ x,\ \Gamma' \vdash M\),\ \ \ \(\Gamma,\ \Gamma' \vdash N\) in \(N \ \#\  M \ \vert\  \Gamma,\ \Gamma'\).

Dokazujemo \(\Gamma,\ \Gamma' \vdash M[N/x]\), torej, da je \(M[N/x]\) izraz v kontekstu \(\Gamma\).

\begin{itemize}
\item Predpostavimo najprej \(M = x\). Potem sledi \(M[N/x] = N\), po definiciji substitucije. Iz \(\Gamma,\ \Gamma' \vdash N\) torej sledi \(M[N/x]\) v kontekstu \(\Gamma,\ \Gamma'\).

\item Predpostavimo zdaj \(M = z\), kjer je \(z \neq x\) spremenljivka. Potem je \(M[N/x] = z\), po definiciji substitucije. Iz predpostavke \(\Gamma,\ x,\ \Gamma' \vdash M\) sledi \(z \in \Gamma,\ \Gamma'\) in torej \(\Gamma,\ \Gamma' \vdash z\).

\item Predpostavimo \(M = M' M''\). Iz predpostavke \(\Gamma,\ x,\ \Gamma' \vdash M\) sledi \(\Gamma,\ x,\ \Gamma' \vdash M'\) in \(\Gamma,\ x,\ \Gamma' \vdash M''\). Po indukcijski predpostavki velja \(\Gamma,\ \Gamma' \vdash M'[N/x]\) in \(\Gamma,\ \Gamma' \vdash M''[N/x]\). Sledi \(\Gamma,\ \Gamma' \vdash M'[N/x] M''[N/x]\) po pravilih za generiranje izrazov. Po definiciji substitucije sledi \(\Gamma,\ \Gamma' \vdash (M' M'')[N/x]\), torej smo dokazali \(\Gamma,\ \Gamma' \vdash M[N/x]\).

\item Predpostavimo \(M = \lambda z. M'\). Po predpostavki velja \(\Gamma,\ x,\ \Gamma' \vdash \lambda z. M'\). Velja \(z \notin \Gamma,\ x,\ \Gamma'\). Po pravilu za generiranje \(\lambda\)-izrazov sledi \(\Gamma,\ x,\ \Gamma',\ z \vdash M'\).

Najprej izpeljimo pogoj \(z \ \#\ \Gamma,\ \Gamma',\ N\).

Velja \(M \ \#\  N \ \vert\  \Gamma,\ \Gamma'\), torej \(\lambda z. M' \ \#\  N \ \vert\  \Gamma,\ \Gamma'\), torej \(z,\ M' \ \#\  N \ \vert\  \Gamma,\ \Gamma'\).

Sledi \(z \ \#\  N \ \vert\  \Gamma,\ \Gamma'\), po pravilu za manipuliranje relacije separiranosti \(\ref{a2}\). Sledi \(z \ \#\  N,\ \Gamma,\ \Gamma'\) (po \ref{a3}), kar smo želeli dokazati (operacija \(\_,\_\) je komutativna).

Podobno lahko iz \(z,\ M' \ \#\  N \ \vert\  \Gamma,\ \Gamma'\) izpeljemo \(M' \ \#\  N \ \vert\  \Gamma,\ \Gamma',\ z\) (po \(\ref{a1}\)).

Ker velja \(z \ \#\  \Gamma,\ \Gamma',\ N\), lahko na predpostavki \(\Gamma,\ \Gamma' \vdash N\) uporabimo lemo, da dobimo \(\Gamma,\ \Gamma',\ z \vdash N\).

Z uporabo dejstev \(\Gamma,\ x,\ \Gamma',\ z \vdash M'\),\ \ \  \(\Gamma,\ \Gamma',\ z \vdash N\),\ \ \ \(M' \ \#\  N \ \vert\  \Gamma,\ \Gamma',\ z\) sledi po indukcijski predpostavki \(\Gamma,\ \Gamma',\ z \vdash M'[N/x]\). Torej sledi \(\Gamma,\ \Gamma' \vdash \lambda z. (M'[N/x])\) in \(\Gamma,\ \Gamma' \vdash M[N/x]\). \qed

\end{itemize}

\begin{opomba}
    \(x \ \#\  y \ \vert\  z\), kjer \(z\) prazno zaporedje, je ekvivalentno \(x \ \#\  y\).
\end{opomba}

\begin{opomba}
    Ker so atomi v zaporedjih, ki predstavljajo kontekst za \(\lambda\)-izraze, enolični (se ne ponavljajo), nesmiselnih izrazov ne moremo skonstruirati, npr. \(\lambda x.(\lambda x.x)\). Lahko pa definiramo izraze kot so \((\lambda x.x)(\lambda x.x)\).
\end{opomba}

Sedaj lahko definiramo še relacijo, ki posploši substitucijo:

\begin{definicija}
    Naj bodo \(\Gamma,\ x,\ \Gamma' \vdash M\),\ \ \ \(\Gamma,\ \Gamma' \vdash N\) in \(\Gamma,\ \Gamma' \vdash L\) lambda izrazi.
    Definirajmo relacijo \(\sub\) med temi izrazi kot:
    \[\sub(M, N, L) \coloneq \exists\ \Gamma,\ \Gamma' \vdash N'.\ \ \ N' \sim N \ \vert\ \Gamma,\ \Gamma',\ \ \ M \ \#\ N' \ \vert \ \Gamma,\ \Gamma',\ \ \ L = M[N'/x].\]
\end{definicija}
Če ta relacija velja, potem označimo
\[
    \Gamma,\ \Gamma' \vdash L \approx M[N/x]
\]

\subsection{\(\alpha\)-ekvivalenca}

Zanimalo nas bo, kdaj sta dva izraza ekvivalentna v nekem smislu: želimo, uvesti ekvivalenco izrazov, ki pove, kdaj sta izraza enaka po pomenu. Natančneje, želimo, da sta dva izraza ekvivalentna, če lahko enega dobimo iz drugega le s preimenovanjem vezanih spremenljivk. Tej ekvivalenci pravimo \(\alpha\)-ekvivalenca.

\begin{definicija}
    Relacija \(\equiv_\alpha\) je definirana s sledečimi pravili:
    \[
        \inference[]{}{\Gamma,\ x,\ \Gamma' \vdash x \equiv_\alpha x}
    \]
    \[
        \inference{\Gamma \vdash M \equiv_\alpha M' \quad \Gamma \vdash N \equiv_\alpha N'}{\Gamma \vdash MN \equiv_\alpha M'N'}
    \]
    \[
        \inference{\Gamma,\ x \vdash M \equiv_\alpha N}{\Gamma \vdash \lambda x.M \equiv_\alpha \lambda x.N}
    \]
    \[
        \inference{\Gamma \vdash M \quad \Gamma \vdash N \quad M \sim N \ \vert\ \Gamma}{\Gamma \vdash M \equiv_\alpha N}
    \]
    Če za dva izraza velja \(\Gamma \vdash M \equiv_\alpha N\), potem rečemo, da sta izraza \textbf{\(\alpha\)-ekvivalentna}.
\end{definicija}
Opomba: \(\equiv_\alpha\) je kontekstualizirana kongruenčna relacija.

Opomba: iz produkcijskih pravil je razvidno, da sta dva izraza lahko \(\alpha\)-ekvivalentna le, če ju opazujemo nad istim kontekstom.

Opomba: relacija \(\equiv_\alpha\) vsebuje kontekstualizirano orbitno ekvivalenco (tj. relacijo \(\sim\)).

Pri zadnjem pravilu velja omeniti naslednjo lastnost.

\begin{trditev}
    Naj \(\Gamma \vdash M\) izraz nad kontekstom \(\Gamma\) in naj bo še \(N\) en drug izraz, tako da velja \(M \sim N \ \vert\ \Gamma\). Potem velja \(\Gamma \vdash N\), tj. izraz \(N\) živi v kontekstu \(\Gamma\).
\end{trditev}

\noindent
{\em Dokaz:\/}
Pogoj \(M \sim N \ \vert\ \Gamma\) implicira obstoj permutacije \(\pi\), ki fiksira člene zaporedja \(\Gamma\) in da velja \(\pi(M) = N\). Sledi \(\pi(\Gamma \vdash M) = \pi(\Gamma) \vdash \pi(M) = \Gamma \vdash N\).
\qed

Opomba pravi, da če permutiramo samo vezane spremenljivke, se kontekst ne spremeni. To je seveda povsem smiselno, namreč kontekst sestavljajo le proste spremenljivke, ne vezane.

\subsection{\(\beta\)-redukcija}

V tem podpoglavju definiramo relacijo \(\beta\)-redukcije. To je funkcijska relacija, ki poveže lambda izraz z lambda izrazom, v katerega se prvotni izraz poenostavi.

\begin{definicija}
    Relacija \(\beta\)-redukcija (označimo z \(\rightarrow_\beta\)) je definirana s sledečimi pravili:
    \[
        \inference{\Gamma \vdash M \rightarrow_\beta M' \quad \Gamma \vdash N}{\Gamma \vdash MN \rightarrow_\beta M'N}
    \]
    \[
        \inference{\Gamma \vdash N \rightarrow_\beta N' \quad \Gamma \vdash M}{\Gamma \vdash MN \rightarrow_\beta MN'}
    \]
    \[
        \inference{\Gamma,\ x \vdash M \rightarrow_\beta M'}{\Gamma \vdash \lambda x.M \rightarrow_\beta \lambda x.M'}
    \]
    \[
        \inference{\Gamma,\ x \vdash M \quad \Gamma \vdash N \quad N' \sim N \ \vert\ \Gamma \quad M \ \#\ N' \ \vert\ \Gamma}{\Gamma \vdash (\lambda x.M)N \rightarrow_\beta M[N'/x]}
    \]
\end{definicija}

Ta zadnje pravilo lahko krajše zapišemo kot
\[
    \inference{\Gamma \vdash L \approx M[N/x]}{\Gamma \vdash (\lambda x.M)N \rightarrow_\beta M[N'/x]}
\]

Naslednja lema bo osrednjega pomena.

\begin{lema}
    Naj \(\Gamma \vdash M \rightarrow_\beta M'\) in naj
    \(\Gamma \vdash M \equiv_\alpha N\).
    Potem velja, da obstaja \(N'\), da \(\Gamma \vdash M' \equiv_\alpha N' \) in \(\Gamma \vdash N \rightarrow_\beta N'\).
\end{lema}

\newpage

\section{Ohranjanje \(\alpha\)-ekvivalence pri substituciji}

Dokazali bomo nekaj pomožnih lastnosti o relaciji \(\alpha\)-ekvivalence, ko v \(\lambda\)-izrazih substituiramo spremenljivke z novim \(\lambda\)-izrazom.

\begin{lema}
    Naj bo \(\Gamma\) kontekst, \(x \in \mathds{A}\) pa atom, ki se ne pojavi v \(\Gamma\). Naj velja \(\Gamma,\ x \vdash M\), \(\Gamma \vdash N,\ N'\), tako da \(\Gamma \vdash N \equiv_\alpha N'\). Naj velja še \(M \ \#\ N \ \vert\ \Gamma\) in \(M \ \#\ N' \ \vert\ \Gamma\). Potem sta \(\Gamma \vdash M[N/x],\ M[N'/x]\) dobro definirana in velja \(\Gamma \vdash M[N/x] \equiv_\alpha M[N'/x]\).
\end{lema}

\noindent
{\em Dokaz:\/}

Dobra definiranost sledi iz definicije substitucije,
\(\alpha\)-ekvivalenco pa lahko dokažemo z indukcijo na \(\lambda\)-izraz \(M\), kot sledi.

\begin{itemize}
    \item Recimo, da \(M = z\), kjer je \(z\) neka spremenljivka. 
        \begin{itemize}
            \item Naj \(x = z\). Potem \(M[N/x] = N\) in \(M[N'/x] = N'\), torej sta po predpostavki \(\alpha\)-ekvivalentna.
            \item Naj \(x \neq z\). Potem \(M[N/x] = z = M[N'/x]\), torej sta \(\alpha\)-ekvivalentna.
        \end{itemize}

    \item Recimo, da \(M = M'M''\), kjer \(\Gamma \vdash M',\ M''\). Velja \(M[N/x] = (M'M'')[N/x] = M'[N/x]M''[N/x]\) in podobno \(M[N'/x] = M'[N'/x]M''[N'/x]\). Po indukcijski predpostavki velja \(M'[N/x] \equiv_\alpha M'[N'/x]\) in \(M''[N/x] \equiv_\alpha M''[N'/x]\), zato po definiciji \(\alpha\)-ekvivalence sledi \(M'[N/x]M''[N/x] \equiv_\alpha M'[N'/x]M''[N'/x]\). Torej imamo \(\alpha\)-ekvivalenco.

    \item Recimo, da \(M = \lambda z. M'\). Velja \(z \neq x\), saj sicer \(\Gamma,\ x \vdash M\) ni definiran. Po indukcijski predpostavki velja \(\Gamma,\ z \vdash M'[N/x] \equiv_\alpha M'[N'/x]\). Po definiciji \(\alpha\)-ekvivalence velja \(\Gamma \vdash \lambda z.M'[N/x] \equiv_\alpha \lambda z.M'[N'/x]\), kar smo želeli dokazati.
\end{itemize}
\qed

Podobna, a manj trivialna je naslednja lema.

\begin{lema}
    Naj bo \(\Gamma\) kontekst, \(x \in \mathds{A}\) pa atom, ki se ne pojavi v \(\Gamma\). Naj velja \(\Gamma,\ x \vdash M,\ M'\), \(\Gamma \vdash N\), tako da \(\Gamma,\ x \vdash M \equiv_\alpha M'\). Naj velja še \(M \ \#\ N \ \vert\ \Gamma\) in \(M' \ \#\ N \ \vert\ \Gamma\). Potem sta \(\Gamma,\ x \vdash M[N/x],\ M'[N/x]\) dobro definirana in velja \(\Gamma \vdash M[N/x] \equiv_\alpha M'[N/x]\).
\end{lema}

Preden dokažemo lemo, uvedimo naslednji koncept:

\begin{definicija}
    Naj funkcija \(f : A \rightarrow B\) slika iz nominalne množice \(A\) v nominalno množico \(B\). Naj \(x,\ y \in A\) in \(\Gamma\) končno zaporedje elementov iz množice atomov, \(\mathds{A}\). Če iz \(x \sim y \ \vert\ \Gamma\), kjer z \(\sim\) označimo relacijo orbitne ekvivalence, sledi \(f(x) \sim f(y) \ \vert\ \Gamma\), potem označimo
    \[
        f \sim f \ \vert\ \Gamma.
    \]
    Podobno, če iz \(x \ \#\ y \ \vert\ \Gamma\) (\(\supp(x)\, \cap\, \supp(y) \subseteq \Gamma\)) sledi \(f(x) \ \#\ f(y) \ \vert\ \Gamma\), potem označimo
    \[
        f \ \#\ f \ \vert\ \Gamma.
    \]
\end{definicija}

Naj bo \(\Gamma\) zaporedje atomov in \(x \in \mathds{A}\) atom, \(x \notin \Gamma\).
Naj bodo \(\Gamma,\ x \vdash M,\ M'\), \(\Gamma \vdash N\), \(M \ \#\ N \ \vert\ \Gamma\) in \(M' \ \#\ N \ \vert\ \Gamma\). Definirajmo funkcijo \(\sub_{x,\ N}\): \[\sub_{x,\ N}(M) \coloneq M[N/x].\]

Če dokažemo \(\sub_{x,\ N} \sim \sub_{x,\ N} \ \vert\ \Gamma,\ x,\ N\) in \(\sub_{x,\ N} \ \#\ \sub_{x,\ N} \ \vert\ \Gamma,\ x,\ N\), potem lahko sklepamo, da iz \(M \sim M' \ \vert\ \Gamma,\ x,\ N\) sledi \(\sub_{x,\ N}(M) \sim \sub_{x,\ N}(M') \ \vert\ \Gamma,\ x,\ N\).


Po definiciji \(\alpha\)-ekvivalence velja, da iz \(\Gamma,\ x \vdash M \equiv_\alpha M'\) sledi \(\Gamma \vdash M[N/x] \equiv_\alpha M'[N/x]\).

Zdaj lahko to lastnost uporabimo pri dokazovanju prejšnje leme.

\noindent
{\em Dokaz:\/}
Če predpostavimo pogoje leme, potem sta izraza \(\Gamma,\ x \vdash M[N/x],\ M'[N/x]\) zares dobro definirana in velja \(\Gamma \vdash M[N/x] \equiv_\alpha M'[N/x]\). 
\qed

Opremljeni s tema dvema lemama se lahko lotimo dokazovanja leme, ki posploši obe prejšnji.

\begin{lema}
    Naj bo \(\Gamma\) kontekst, \(x \in \mathds{A}\) pa atom, ki se ne pojavi v \(\Gamma\). Naj velja \(\Gamma,\ x \vdash M,\ M'\) in \(\Gamma \vdash N,\ N'\), tako da \(\Gamma \vdash M \equiv_\alpha M'\) in \(\Gamma \vdash N \equiv_\alpha N'\). Naj velja še \(M \ \#\ N \ \vert\ \Gamma\) in \(M' \ \#\ N' \ \vert\ \Gamma\). Potem sta \(\Gamma \vdash M[N/x],\ M'[N'/x]\) dobro definirana in velja \(\Gamma \vdash M[N/x] \equiv_\alpha M'[N'/x]\).
\end{lema}

\noindent
{\em Dokaz:\/}

Naj veljajo predpostavke iz formulacije leme.
Ideja je, da uporabimo ta drugo lemo, da zamenjamo \(M\) v \(M'\), substituenta pa zamenjamo s ta prvo lemo. Izkaže se, da bomo prvo lemo uporabili dvakrat, da izpolnimo pogoje za definiranost substitucije.

Uporabimo trditev \ref{pogojna_trditev} na nominalni množici \(\Terms_\Gamma\), kjer za \(x\) vzamemo \(N\), za \(y\) vzamemo \((M,\ M')\) in za \(w\) vzamemo \(\Gamma\). Iz trditve potem sledi, da obstaja \(L\), tako da \(L \sim N \ \vert\ \Gamma\) in \(L \ \#\ M,\ M' \ \vert\ \Gamma\).

Iz definicije \(\alpha\)-ekvivalence sledi, da \(L \equiv_\alpha N\). Ker predpostavljamo \(N \equiv_\alpha N'\), sledi \(L \equiv_\alpha N'\), po tranzitivnosti \(\alpha\)-ekvivalence.

Sedaj lahko sklepamo:
\[
    \Gamma \vdash M[N/x] \equiv_\alpha M[L/x] \equiv_\alpha M'[L/x] \equiv_\alpha M'[N'/x],
\]
kjer smo pri prvi in tretji ekvivalenci uporabili prvo lemo z dejstvom, da \(N \equiv_\alpha L \equiv_\alpha N'\), pri drugi ekvivalenci pa smo uporabili drugo lemo.

\qed
\newpage
\section{Church-Rosserjev izrek}
Preden formuliramo ciljni izrek, definirajmo naslednje:
\begin{oznaka}
    Refleksivno-tranzitivno ovojnico relacije \(\rightarrow_\beta\) označimo z \(\rightarrow_\beta^*\). To je torej najmanjša relacija \(\rightarrow_\beta^*\), da za vsak kontekst \(\Gamma\) in izraze \(\Gamma \vdash M,\ N,\ L\) veljajo
    \[
        \Gamma \vdash M \rightarrow_\beta N\, \Longrightarrow\, \Gamma \vdash M \rightarrow_\beta^* N,
    \]
    \[
        \Gamma \vdash M \rightarrow_\beta^* M,
    \]
    \[
        \Gamma \vdash M \rightarrow_\beta^* N\, \wedge\, \Gamma \vdash N \rightarrow_\beta^* L\, \Longrightarrow\, \Gamma \vdash M \rightarrow_\beta^* L.
    \]
\end{oznaka}

Zdaj lahko formuliramo glavni izrek.
\begin{izrek_crosser}
    Naj bo \(\Gamma\) kontekst in \(\Gamma \vdash M,\ N_1,\ N_2\) \(\lambda\)-izrazi nad tem kontekstom. Naj \(\Gamma \vdash M \rightarrow_\beta^* N_1\) in \(\Gamma \vdash M \rightarrow_\beta^* N_2\). Potem obstaja \(\Gamma \vdash X\), da velja
    \[
        \Gamma \vdash N_1 \rightarrow_\beta^* X
    \]
    in
    \[
        \Gamma \vdash N_2 \rightarrow_\beta^* X.
    \]
\end{izrek_crosser}


\noindent
{\em Dokaz:\/}
Dokaz zaenkrat opustimo.
\qed


\end{document}

